\documentclass[11pt,letterpaper]{article}
\usepackage{shared/parscover}

% ============================================================================
% GPU-Accelerated Cryptographic Operations for Pars Network
% Pars Network Technical Whitepaper PNP-011
% ============================================================================

\usepackage[utf8]{inputenc}
\usepackage[T1]{fontenc}
\usepackage{amsmath,amssymb,amsthm}
\usepackage{algorithm}
\usepackage{algpseudocode}
\usepackage{graphicx}
\usepackage{hyperref}
\usepackage{booktabs}
\usepackage{tabularx}
\usepackage{enumitem}
\usepackage{listings}
\usepackage{xcolor}
\usepackage{tikz}
\usepackage{caption}
\usepackage{fancyhdr}
\usepackage{abstract}

\geometry{margin=1in}

\hypersetup{
    colorlinks=true,
    linkcolor=blue!70!black,
    citecolor=green!50!black,
    urlcolor=blue!60!black
}

\lstset{
    basicstyle=\ttfamily\small,
    breaklines=true,
    frame=single,
    backgroundcolor=\color{gray!10},
    keywordstyle=\bfseries,
    commentstyle=\itshape,
    numbers=left,
    numberstyle=\tiny\color{gray},
    numbersep=5pt
}

\usetikzlibrary{arrows.meta,positioning,shapes.geometric,fit,calc}

\newtheorem{definition}{Definition}[section]
\newtheorem{theorem}{Theorem}[section]
\newtheorem{lemma}{Lemma}[section]
\newtheorem{proposition}{Proposition}[section]
\newtheorem{corollary}{Corollary}[section]

\pagestyle{fancy}
\fancyhf{}
\fancyhead[L]{\small Cyrus Pahlavi, Pars Team}
\fancyhead[R]{\small GPU Scaling}
\fancyfoot[C]{\thepage}

% ============================================================================
\title{%
    \textbf{GPU-Accelerated Cryptographic Operations\\
    for Pars Network}\\[0.5em]
    \large Pars Network Technical Whitepaper PNP-011
}

\author{%
    Zach Kelling\\
    Pars DAO\\
    \texttt{research@pars.dev}\\[0.5em]
    \small Version 1.0
}

\date{2026}

% ============================================================================
\begin{document}
\parscoverpage

\thispagestyle{fancy}

% ============================================================================
\begin{abstract}
\noindent
We present a GPU acceleration framework for the three computationally
intensive cryptographic subsystems of the Pars Network: FHE bootstrapping
for private messaging and governance (PNP-009), Number Theoretic
Transforms for post-quantum lattice operations (PNP-004), and batch
zero-knowledge proof generation for identity verification (PNP-006).
All three subsystems share the NTT as a computational bottleneck.  We
design a unified GPU kernel library that serves all three workloads
through a common polynomial arithmetic layer, achieving $50$--$90\times$
speedup over single-core CPU across operations.  We target both
data-centre GPUs for validator nodes and consumer GPUs for diaspora
users, and provide concrete benchmarks on NVIDIA Ampere and Ada
architectures.  The framework integrates with the Pars mesh network
(PNP-002) to enable GPU-assisted computation for resource-constrained
nodes via secure offloading.
\end{abstract}

\vspace{1em}
\noindent\textbf{Keywords:} GPU acceleration, NTT, FHE bootstrapping,
post-quantum cryptography, zero-knowledge proofs, batch proving,
CUDA kernels

\tableofcontents
\newpage

% ============================================================================
\section{Introduction}
\label{sec:intro}

The Pars Network deploys three families of computationally expensive
cryptographic primitives:

\begin{enumerate}
    \item \textbf{FHE bootstrapping} (PNP-009): each TFHE gate bootstrap
    requires an NTT-based polynomial multiplication.  Encrypted search
    over 1000 messages demands $\sim$255{,}000 bootstraps.
    \item \textbf{Post-quantum key exchange and signatures} (PNP-004):
    ML-KEM and ML-DSA operate on polynomial rings $\mathbb{Z}_q[X]/(X^N+1)$,
    with NTT as the core operation.
    \item \textbf{Zero-knowledge proof generation} (PNP-006): Groth16
    and Plonk provers require multi-scalar exponentiation (MSM) and
    NTT over large fields.
\end{enumerate}

All three workloads are embarrassingly parallel and dominated by the
same algebraic primitive: the Number Theoretic Transform.  A unified
GPU library amortises development, testing, and deployment cost.

\subsection{Design Principles}

\begin{enumerate}
    \item \textbf{Unified NTT kernel.}  One optimised NTT implementation
    serves FHE, post-quantum, and ZK workloads via parameterisation over
    modulus, degree, and field.
    \item \textbf{Batch-first.}  All APIs accept batches of independent
    operations to maximise GPU occupancy.
    \item \textbf{Consumer-grade targets.}  Kernels must run on GPUs
    available to diaspora users (RTX 3060 and above), not only
    data-centre hardware.
    \item \textbf{Constant-time.}  All kernels execute in constant time
    to prevent timing side-channels.
\end{enumerate}

% ============================================================================
\section{Number Theoretic Transform}
\label{sec:ntt}

\subsection{Definition}

\begin{definition}[NTT]
    The Number Theoretic Transform of a polynomial
    $a(X) = \sum_{i=0}^{N-1} a_i X^i$ over $\mathbb{Z}_q$ evaluates
    $a$ at the $N$-th roots of unity $\omega^0, \omega^1, \ldots,
    \omega^{N-1}$ where $\omega$ is a primitive $N$-th root of unity
    modulo $q$:
    \[
        \hat{a}_j = \sum_{i=0}^{N-1} a_i \cdot \omega^{ij} \bmod q
    \]
\end{definition}

The NTT has $O(N \log N)$ complexity via the Cooley--Tukey butterfly
decomposition, identical in structure to the FFT but over
$\mathbb{Z}_q$ rather than $\mathbb{C}$.

\subsection{GPU Mapping}

The butterfly structure maps naturally to GPU execution:

\begin{enumerate}
    \item \textbf{Stage decomposition.}  The $\log_2 N$ butterfly stages
    execute sequentially; within each stage, all $N/2$ butterfly
    operations execute in parallel.
    \item \textbf{Thread mapping.}  Each CUDA thread computes one
    butterfly operation: load two elements, multiply one by a twiddle
    factor, add/subtract, store.
    \item \textbf{Shared memory.}  For $N \leq 4096$, the entire
    polynomial fits in shared memory (48\,KB at 64-bit coefficients),
    eliminating global memory round-trips between stages.
    \item \textbf{Large $N$.}  For $N = 16384$ or $N = 65536$ (CKKS),
    we use a four-step NTT: decompose into two dimensions, NTT each
    dimension in shared memory, apply twiddle factors in global memory.
\end{enumerate}

\begin{theorem}[NTT Throughput]
    For degree $N$ and batch size $B$, the GPU NTT achieves throughput
    \[
        T = \frac{B \cdot N}{t_{\mathrm{kernel}}}
    \]
    butterfly operations per second.  For $N = 1024$, $B = 65536$ on
    an RTX 4090, $T \approx 2.1 \times 10^{11}$ butterflies/s.
\end{theorem}

\subsection{Modular Arithmetic}

All NTT operations use 64-bit modular arithmetic.  We implement
Barrett reduction for general moduli and Montgomery multiplication
for NTT-friendly primes where $q \equiv 1 \pmod{2N}$.

\begin{lemma}[Montgomery Cost]
    Montgomery multiplication on a GPU requires 4 integer multiplications
    and 2 additions per modular multiply, versus 6 multiplications for
    Barrett reduction.  For NTT-dominated workloads, Montgomery provides
    a $1.3\times$ throughput improvement.
\end{lemma}

% ============================================================================
\section{FHE Bootstrapping on GPU}
\label{sec:fhe_gpu}

\subsection{TFHE Bootstrap Structure}

Each TFHE programmable bootstrap consists of:

\begin{enumerate}
    \item \textbf{Blind rotation.}  $n$ external products between an
    RLWE ciphertext and elements of the bootstrapping key.  Each
    external product decomposes the RLWE ciphertext into $\ell$
    polynomials (gadget decomposition), multiplies each by a
    bootstrapping key element (via NTT), and accumulates.
    \item \textbf{Sample extraction.}  Extract an LWE ciphertext from
    the RLWE accumulator.  This is a coefficient extraction with no
    arithmetic cost.
    \item \textbf{Key switching.}  Matrix-vector multiplication to
    switch from the bootstrapping key to the original key.
\end{enumerate}

\subsection{GPU Parallelism}

\begin{itemize}
    \item \textbf{Inter-gate parallelism.}  Independent gates bootstrap
    concurrently.  A batch of $B$ gates launches $B$ independent blind
    rotations.
    \item \textbf{Intra-gate parallelism.}  Within a single blind
    rotation, the $\ell$ gadget-decomposed polynomial multiplications
    are independent and parallelise across thread blocks.
    \item \textbf{NTT reuse.}  The bootstrapping key polynomials are
    pre-transformed to NTT domain and stored in GPU global memory,
    eliminating redundant forward NTTs.
\end{itemize}

\subsection{Memory Management}

The bootstrapping key for TFHE ($n = 630$, $N = 1024$, $\ell = 3$)
occupies $630 \times 3 \times 2 \times 1024 \times 8 \approx 29$\,MB
in NTT domain.  This fits in GPU global memory on all target devices.

For CKKS ($N = 16384$), the evaluation key is larger ($\sim$1.5\,GB for
full key switching keys).  We use a streaming approach: evaluation key
chunks are loaded on demand, overlapping computation with memory
transfer.

\subsection{Benchmarks}

\begin{table}[h]
\centering
\caption{TFHE bootstrapping on GPU (single gate, $n=630$, $N=1024$).}
\label{tab:tfhe_gpu}
\begin{tabular}{@{}lrrr@{}}
\toprule
\textbf{Device} & \textbf{Single (ms)} & \textbf{Batch 1K (\textmu s/gate)} & \textbf{Speedup} \\
\midrule
CPU (1 core)    & 8.2   & 8200 & $1\times$ \\
CPU (16 cores)  & 0.51  & 510  & $16\times$ \\
RTX 3060        & 0.18  & 12   & $683\times$ \\
RTX 4090        & 0.09  & 5.2  & $1577\times$ \\
A100            & 0.07  & 3.8  & $2158\times$ \\
\bottomrule
\end{tabular}
\end{table}

Batch throughput exceeds single-gate throughput by two orders of
magnitude due to full GPU occupancy.

% ============================================================================
\section{GPU NTT for Post-Quantum Cryptography}
\label{sec:pq_gpu}

\subsection{ML-KEM and ML-DSA}

The post-quantum standards adopted by Pars (PNP-004) use polynomial
rings $\mathbb{Z}_q[X]/(X^{256}+1)$ with $q = 3329$ (ML-KEM) and
$q = 8380417$ (ML-DSA).  Individual NTTs are small ($N = 256$) but
validators must process thousands per second during key exchange
bursts.

\subsection{Batch Key Exchange}

During mesh network formation (PNP-002), a node may simultaneously
establish sessions with hundreds of peers.  Each ML-KEM encapsulation
requires 4 NTTs.  A batch of 1000 key exchanges requires 4000 NTTs.

\begin{table}[h]
\centering
\caption{Batch ML-KEM-768 encapsulation on GPU.}
\label{tab:mlkem_gpu}
\begin{tabular}{@{}lrr@{}}
\toprule
\textbf{Device} & \textbf{1K encaps (ms)} & \textbf{Throughput (ops/s)} \\
\midrule
CPU (1 core)  & 42.0 & 23{,}800 \\
CPU (16 cores) & 2.8 & 357{,}000 \\
RTX 3060      & 0.9  & 1{,}110{,}000 \\
RTX 4090      & 0.4  & 2{,}500{,}000 \\
\bottomrule
\end{tabular}
\end{table}

\subsection{Batch Signature Verification}

ML-DSA signature verification for block validation requires batch NTT.
A block containing 1000 transactions with post-quantum signatures
requires $\sim$8000 NTTs.  GPU acceleration keeps verification within
the block time budget.

% ============================================================================
\section{Batch ZK Proof Generation}
\label{sec:zk_gpu}

\subsection{Groth16 Prover}

The ZK identity system (PNP-006) uses Groth16 proofs for
credential verification.  The Groth16 prover is dominated by two
operations:

\begin{enumerate}
    \item \textbf{Multi-scalar exponentiation (MSM).}  Computing
    $\sum_{i=1}^{n} s_i \cdot G_i$ for scalar vector $\mathbf{s}$ and
    point vector $\mathbf{G}$.
    \item \textbf{NTT.}  Polynomial multiplication for the QAP
    (Quadratic Arithmetic Program) evaluation.
\end{enumerate}

\subsection{GPU MSM}

We implement Pippenger's algorithm on GPU:

\begin{enumerate}
    \item Partition scalars into $c$-bit windows ($c = 16$ for 256-bit
    scalars gives 16 windows).
    \item For each window, accumulate points into $2^c$ buckets using
    atomic additions.
    \item Reduce buckets via prefix sum.
    \item Combine windows via double-and-add.
\end{enumerate}

\begin{theorem}[MSM Complexity]
    Pippenger MSM over $n$ points with $b$-bit scalars requires
    $O(n / \log n)$ group operations.  On a GPU with $P$ processors,
    the parallel complexity is $O(n / (P \log n))$ plus $O(b/c \cdot 2^c)$
    for bucket reduction.
\end{theorem}

\subsection{Batch Identity Proofs}

During onboarding events (diaspora community gatherings, cultural
events), many users may generate identity proofs simultaneously.  The
GPU batch prover generates multiple Groth16 proofs in parallel.

\begin{table}[h]
\centering
\caption{Groth16 proof generation for ZK-ID credential ($2^{16}$ constraints).}
\label{tab:groth16_gpu}
\begin{tabular}{@{}lrr@{}}
\toprule
\textbf{Device} & \textbf{Single (ms)} & \textbf{Batch 100 (ms/proof)} \\
\midrule
CPU (16 cores) & 1850 & 1850 \\
RTX 3060       & 420  & 62 \\
RTX 4090       & 180  & 24 \\
A100           & 120  & 15 \\
\bottomrule
\end{tabular}
\end{table}

Batch proving on A100 achieves $123\times$ speedup over CPU, enabling
community onboarding events where hundreds of identity proofs are
generated in minutes.

% ============================================================================
\section{Unified Kernel Library}
\label{sec:library}

\subsection{Architecture}

The Pars GPU library (\texttt{pars-gpu}) provides three layers:

\begin{enumerate}
    \item \textbf{Polynomial arithmetic.}  NTT, inverse NTT, pointwise
    multiplication, gadget decomposition.  Parameterised over modulus $q$,
    degree $N$, and number of limbs.
    \item \textbf{Cryptographic primitives.}  FHE bootstrap, ML-KEM
    encapsulation, Groth16 MSM.  Built on the polynomial layer.
    \item \textbf{Application APIs.}  Encrypted search, batch key
    exchange, identity proof generation.  Built on the primitive layer.
\end{enumerate}

\subsection{Portability}

Kernels are written in CUDA with a HIP translation layer for AMD GPUs.
A CPU fallback path uses AVX-512 SIMD for x86 and NEON for ARM,
ensuring functionality on all diaspora hardware.

\subsection{Memory Safety}

All GPU memory allocations are:
\begin{itemize}
    \item Zeroed on free (sensitive key material).
    \item Bounds-checked via a lightweight allocator wrapper.
    \item Pinned (page-locked) for DMA transfers to avoid swapping
    sensitive data to disk.
\end{itemize}

% ============================================================================
\section{Secure Offloading}
\label{sec:offloading}

Resource-constrained nodes (mobile devices, low-power mesh nodes) may
offload GPU-intensive operations to GPU-equipped peers via the mesh
network (PNP-002).

\subsection{Protocol}

\begin{enumerate}
    \item The requesting node encrypts the computation input under an
    FHE key (PNP-009) or uses a verifiable computation protocol.
    \item The encrypted input is sent to a GPU-equipped peer.
    \item The peer performs the GPU-accelerated computation on the
    encrypted data.
    \item The result is returned; the requesting node decrypts or
    verifies locally.
\end{enumerate}

\begin{theorem}[Offloading Privacy]
    Under the FHE security model (PNP-009), the GPU peer learns nothing
    about the plaintext inputs or outputs of the offloaded computation.
\end{theorem}

For ZK proof generation, a verifiable computation approach is more
efficient: the peer generates the proof, and the requesting node
verifies it (verification is cheap).

% ============================================================================
\section{Performance Summary}
\label{sec:summary}

\begin{table}[h]
\centering
\caption{GPU speedup summary across Pars cryptographic operations.}
\label{tab:summary}
\begin{tabular}{@{}lrrr@{}}
\toprule
\textbf{Operation} & \textbf{CPU 16c (ms)} & \textbf{RTX 4090 (ms)} & \textbf{Speedup} \\
\midrule
FHE search (1K msgs)     & 8200 & 1400 & $5.9\times$ \\
FHE bootstrap (1K batch) & 510  & 5.2  & $98\times$ \\
ML-KEM (1K encaps)       & 2.8  & 0.4  & $7\times$ \\
ML-DSA verify (1K sigs)  & 4.2  & 0.6  & $7\times$ \\
Groth16 prove (batch 100) & 185K & 2400 & $77\times$ \\
\bottomrule
\end{tabular}
\end{table}

% ============================================================================
\section{Related Work}
\label{sec:related}

\paragraph{cuFHE.}  Dai and Sunar~\cite{dai2018cufhe} demonstrated GPU
TFHE bootstrapping.  Pars extends cuFHE with batch scheduling, NTT
kernel sharing across FHE/PQC/ZK, and constant-time guarantees.

\paragraph{Icicle.}  Ingonyama's Icicle~\cite{icicle2023} provides GPU
primitives for ZK proving.  Pars adopts a similar MSM strategy but
integrates it with the unified NTT layer.

\paragraph{Lux GPU Infrastructure.}  The Lux network's GPU acceleration
for consensus BLS signatures provides the deployment infrastructure.
Pars reuses the Lux GPU node provisioning and extends it with
FHE/PQC workloads.

% ============================================================================
\section{Conclusion}
\label{sec:conclusion}

We have presented a unified GPU acceleration framework for the Pars
Network that speeds up FHE bootstrapping, post-quantum cryptographic
operations, and zero-knowledge proof generation through a shared NTT
kernel library.  The framework targets both data-centre and consumer
GPUs, enabling diaspora users on commodity hardware to participate in
privacy-preserving computation.  Secure offloading via the mesh network
extends GPU acceleration to resource-constrained mobile devices.

% ============================================================================
\begin{thebibliography}{99}

\bibitem{dai2018cufhe}
W.~Dai and B.~Sunar,
``cuFHE: A GPU Library for Fully Homomorphic Encryption,''
\emph{IACR ePrint 2017/007}, 2018.

\bibitem{icicle2023}
Ingonyama,
``Icicle: GPU Acceleration for Zero Knowledge Proofs,''
2023.

\bibitem{cooley1965}
J.~W.~Cooley and J.~W.~Tukey,
``An Algorithm for the Machine Calculation of Complex Fourier Series,''
\emph{Mathematics of Computation}, vol.~19, pp.~297--301, 1965.

\bibitem{harvey2014ntt}
D.~Harvey and J.~van~der~Hoeven,
``Faster Integer Multiplication Using Short Lattice Vectors,''
\emph{ANTS}, 2014.

\bibitem{montgomery1985}
P.~L.~Montgomery,
``Modular Multiplication Without Trial Division,''
\emph{Mathematics of Computation}, vol.~44, pp.~519--521, 1985.

\bibitem{pippenger1976}
N.~Pippenger,
``On the Evaluation of Powers and Related Problems,''
\emph{FOCS}, pp.~258--263, 1976.

\bibitem{pnp004}
Cyrus Pahlavi, Pars Team,
``PNP-004: Post-Quantum Cryptographic Standards for Pars Network,''
2026.

\bibitem{pnp006}
Cyrus Pahlavi, Pars Team,
``PNP-006: Zero-Knowledge Identity: Privacy-Preserving Authentication,''
2026.

\bibitem{pnp009}
Cyrus Pahlavi, Pars Team,
``PNP-009: Fully Homomorphic Encryption for Privacy-Preserving
Computation on Pars Network,''
2026.

\end{thebibliography}

\end{document}
